\section{Matched-Grid Display and Learning Diagnostics}
\label{app:revision-diagnostics}

\subsection{Complete cell evidence and deployment interpretation}
% evidence: paper/generated/reviewer_revision_v12/figure_manifest.json
The display in Figure~\ref{fig:matched-grid} is generated directly from
\texttt{table\_added\_value\_contrasts.csv}. It preserves all 300
task--setting--dose rows, with their paired pointwise CCC/MAE intervals and
within-model dose-family CCC bands. The four resource corners give 48
comparisons across all settings, including unfavorable and uncertain results.
The files \texttt{matched\_increment\_all\_cells.csv} and
\texttt{deployment\_corners.csv} retain unrounded values. A common symmetric
color scale is used, with entries multiplied by 1000 for legibility.
The visualization was added after examining the results and is descriptive.

The four corners answer different deployment questions: no target-user
calibration or repeated-stimulus ratings; calibration only; repeated-stimulus
ratings only; and both. Released normative information remains separately
declared. These scenarios do not make ratings costless or assert that stimuli
will recur in a real application. A weighted application score requires
weights selected for that application and paired prediction-level uncertainty.
Marginal interval endpoints cannot be averaged to obtain such an interval.

\subsection{Real-task optimization evidence}
% evidence: paper/generated/reviewer_revision_v12/learning_diagnostics.csv
\begin{table}[ht]\centering\small
\caption{Previously executed EEGNet training diagnostics, now summarized by
task and branch. Counts refer to settings, identity blocks, and stimulus doses,
not independent participants or studies. Objective reduction is the median
initial-minus-final recorded training objective; these values are not test
metrics and are not comparable performance measures across tasks.}
\begin{tabular}{llrrrr}\toprule
Task & Branch & Fits & Epoch-one & Median epoch & Objective reduction\\\midrule
Emo & Direct & 50 & 12 & 18 & 0.1104\\
Emo & Residual & 50 & 10 & 17.5 & 0.1162\\
Urban & Direct & 50 & 8 & 11 & 0.1922\\
Urban & Residual & 50 & 7 & 7.5 & 0.2080\\
\bottomrule\end{tabular}
\end{table}
These aggregates come from the same verification record as the primary
predictions, not from newly selected runs. The historical standard/compact
EEGNet and matched shrinkage-calibration comparisons remain supporting
evidence on their original supports. The repaired real SEED-V task provides
a separate compatibility check, not a positive control for the current
self-report regression task. None of these records spans full fine-tuning,
all calibration algorithms, or repeated-training uncertainty. The scope of
the matched increment remains the specified fitted predictor pair.

\subsection{Post-primary partition, initialization, and fine-tuning checks}
% evidence: paper/generated/final_revision_v13/sensitivity/sensitivity_summary.json; paper/generated/final_revision_v13/sensitivity/finetune_execution
\begin{table}[ht]\centering\small
\caption{Post-primary robustness checks. Block-$t$ intervals remain conditional on each fitted run. Alternative partitions and initialization schedules are sensitivity analyses, not independent cohorts.}
\begin{tabularx}{\linewidth}{p{0.09\linewidth}p{0.22\linewidth}p{0.18\linewidth}X}\toprule
Task & Check & Setting & Result\\\midrule
Emo & New blocks + refit & Six settings & max point shift 0.0135; 6/6 cover 0\\
Urban & New blocks + refit & Six settings & max point shift 0.0154; 6/6 cover 0\\
Emo & Three train seeds & EEGNet L1 & [-0.0094, -0.0015]; 3/3 cover 0\\
Emo & Three train seeds & EEGNet L2 & [-0.0018, +0.0003]; 3/3 cover 0\\
Urban & Three train seeds & EEGNet L1 & [+0.0097, +0.0134]; 3/3 cover 0\\
Urban & Three train seeds & EEGNet L2 & [+0.0013, +0.0109]; 3/3 cover 0\\
Emo & Last-block fine-tune & LaBraM residual & -0.0110 [-0.1041, +0.0820]\\
Urban & Last-block fine-tune & LaBraM residual & -0.0404 [-0.1159, +0.0351]\\
\bottomrule\end{tabularx}
\end{table}

The alternative partition repeats assignment, fitting, aggregation, and
verification with a new design seed; because the fitted data also change, it is
a partition/refit sensitivity rather than an isolated partition effect. The
three EEGNet schedules share the original assignments and four resource
corners, changing only declared training initialization offsets. The LaBraM
check uses the official checkpoint, freezes all but the final transformer block
and normalization, fits a residual head, selects epochs on validation data, and
then refits on training plus validation. Its 100-Hz analysis tensors are
linearly resampled to the official 200-Hz patch interface; the Urban voltage
conversion is checked from recorded metadata. All checks condition their
block-$t$ intervals on the resulting fitted predictions. They were specified
after the primary results, cover neither a new cohort nor the full $5\times5$
fine-tuning surface, and cannot certify invariance to arbitrary partitions,
optimizers, checkpoints, or adaptation methods.

\subsection{Five-block interval validation}
The production statistic first computes paired CCC differences within each
identity block, then averages blocks and dose cells. It does not pool all
trials before computing CCC. The new validation implementation preserves this
order and uses the same global participant and stimulus multinomial weights
across all blocks and cells. A draw invalid in any block is not repaired by
omitting that block. Tests compare its cell estimates, uniform means, interval
endpoints, and valid-draw counts with the production analyzer for both one-
and two-target tasks.

% evidence: configs/block_validation_v12.json; paper/generated/final_closure_v9/percentile/added_value_statistics.json
The frozen validation configuration uses the exact complete diagonal test
supports in the primary verification record: $6\times9,6\times9,6\times8,
6\times8,6\times8$ for Emo, and $13\times12,13\times11,12\times11,
12\times11,12\times11$ for Urban. It retains the full $5\times5$ dose grid,
two targets for Emo and one for Urban. Four previously defined Gaussian
scenarios exercise label-only cancellation, trial signal, participant-constant
signal, and independent noise. Each scenario uses 300 fresh simulation
replicates per task and 2,000 bootstrap draws. The primary block-$t$ interval
retains the crossed-bootstrap standard error and uses the five identity blocks
as the independent aggregation units (df=4). Percentile, basic, and the earlier
within-block t-normal intervals remain sensitivity checks. The correction was
introduced after the initial interval diagnostics failed, but its critical
value is fixed by the estimand topology and does not use the real EEG effect
signs or magnitudes.

The estimand is the scenario's analytic population matched CCC difference,
averaged uniformly over dose cells. Finite-block estimation bias is part of
the stress test. Synthetic block resources are independent: matching this
test topology does not reproduce all cross-block training-resource dependence,
the actual EEG distribution, or retraining uncertainty. Coverage proportions
require Monte Carlo uncertainty and do not establish validity outside the
tested constructions. Executed results are recorded in the accompanying
validation artifact.


% evidence: results/block_t_validation_v13/summary.csv; results/block_t_validation_v13/completion.json
\begin{table}[ht]\centering\small
\caption{Coverage of the primary uniform-grid mean under the observed five-block topology (300 replicates per row). Brackets give Wilson Monte Carlo intervals for the primary block-$t$ coverage.}
\begin{tabular}{llccc}\toprule
Task & Scenario & Block-$t$ [MC interval] & Percentile & Basic\\\midrule
Emo & Trial signal & 0.993 [0.976, 0.998] & 0.900 & 0.983\\
Emo & Participant constant & 0.980 [0.957, 0.991] & 0.903 & 0.910\\
Emo & Independent noise & 1.000 [0.987, 1.000] & 1.000 & 1.000\\
Emo & Labels only & 1.000 [0.987, 1.000] & 1.000 & 1.000\\
Urban & Trial signal & 1.000 [0.987, 1.000] & 0.927 & 0.977\\
Urban & Participant constant & 0.983 [0.962, 0.993] & 0.970 & 0.920\\
Urban & Independent noise & 1.000 [0.987, 1.000] & 1.000 & 1.000\\
Urban & Labels only & 1.000 [0.987, 1.000] & 1.000 & 1.000\\
\bottomrule\end{tabular}
\end{table}

All 2,400 simulation replicates completed without a failed run. The original
187,200 candidate rows and 62,400 derived block-$t$ rows include every cell and
uniform mean; they are not 249,600 independent simulations. Labels-only cancellation gives degenerate
zero differences; its coverage of one is an algebraic check, not calibration
of a nondegenerate interval. Independent noise produces no positive interval
in these runs, which does not establish false-positive control for every null.

% evidence: results/block_t_validation_v13/summary.csv; paper/generated/reviewer_revision_v12/five_block_all_coverage.csv
The earlier percentile failure is substantive. At Emo's zero-dose
participant-constant cell, its population-target bias is $-0.0756$ and coverage
is $0.193$ (Monte Carlo interval $[0.153,0.242]$); the earlier t-normal coverage
is $0.887$. The block-$t$ construction raises the minimum cell coverage to
$0.917$, which still does not validate cellwise inference. Its uniform-mean
coverage is $0.980$--$1.000$ across the eight task--scenario combinations, with
Wilson lower bounds of $0.957$--$0.987$. These tests support the primary
uniform-grid mean under the simulated topology, not every cell, the actual EEG
distribution, arbitrary signal families, or retraining uncertainty. Resource
identity and descriptive paired differences do not depend on a nominal
significance threshold; no equivalence claim is made.
